%%%%%%%%%%%%%%%%%
% CV tailored for Davidson - Chef de projet SI
%%%%%%%%%%%%%%%%%

\documentclass[8pt,a4paper,ragged2e,withhyper]{altacv}

\geometry{left=1.25cm,right=1.25cm,top=.5cm,bottom=1.cm,columnsep=1.2cm}

\usepackage{paracol}
\usepackage{fontawesome5}
\usepackage{hyperref}

\iftutex 
  \setmainfont{Roboto Slab}
  \setsansfont{Lato}
  \renewcommand{\familydefault}{\sfdefault}
\else
  \usepackage[rm]{roboto}
  \usepackage[defaultsans]{lato}
  \renewcommand{\familydefault}{\sfdefault}
\fi

\definecolor{SlateGrey}{HTML}{2E2E2E}
\definecolor{LightGrey}{HTML}{666666}
\definecolor{DarkPastelRed}{HTML}{8AC0D9}
\definecolor{PastelRed}{HTML}{00B0DA}
\definecolor{GoldenEarth}{HTML}{B4AD0C}

\colorlet{name}{black}
\colorlet{tagline}{PastelRed}
\colorlet{heading}{DarkPastelRed}
\colorlet{headingrule}{GoldenEarth}
\colorlet{subheading}{PastelRed}
\colorlet{accent}{PastelRed}
\colorlet{emphasis}{SlateGrey}
\colorlet{body}{LightGrey}

\definecolor{violet}{HTML}{5A2582}
\definecolor{rouge}{HTML}{F53B3B}
\definecolor{noir}{HTML}{00232C}
\definecolor{bleu}{HTML}{00B0DA}
\definecolor{rose}{HTML}{F831A0}
\definecolor{jaune}{HTML}{FFEC00}
\definecolor{vert}{HTML}{34AD6C}

\renewcommand{\namefont}{\Huge\rmfamily\bfseries}
\renewcommand{\personalinfofont}{\footnotesize}
\renewcommand{\cvsectionfont}{\LARGE\rmfamily\bfseries}
\renewcommand{\cvsubsectionfont}{\large\bfseries}
\renewcommand{\cvItemMarker}{{\small\textbullet}}
\renewcommand{\cvRatingMarker}{\faCircle}

\input{../../_shared/pubs-num.tex}
\addbibresource{../../_shared/sample.bib}

\begin{document}

\name{BOILEAU Guillaume}
\tagline{Chef de Projet SI | Cadrage, Recette, Coordination Technique, Reporting \& Conduite du Changement}

\photoL{2.8cm}{../../_shared/moi}

\personalinfo{%
  \email{guillaume.boileaupro@gmail.com}
  \phone{+33 6 59 94 85 84}
  \location{Alpes-Maritimes, FRANCE}
  \linkedin{guillaume-boileau-phd-891b81265}
  \researchgate{Guillaume-Boileau-2}
  \orcid{0000-0002-3576-6968}
}

\makecvheader
\vspace{-0.3cm}

\AtBeginEnvironment{itemize}{\small}
\columnratio{0.64}

\begin{paracol}{2}

\cvsection{Experience}

\cvevent{\textbf{Software Engineer - System Integration, Delivery \& Production Follow-up}}{VEV Platform Services France}{2025 -- 2026}{Valbonne, France}

\textbf{Context:} Development and integration of software services for an EV charging CPMS platform connected to operational infrastructure and field equipment.

\textbf{Objective:} Support reliable delivery, production follow-up, software integration and operational data-flow continuity between platform services and charging station systems.

\begin{itemize}
  \item \textbf{Software delivery:} Developed and maintained web and backend services using TypeScript, Angular and Node.js in an operational industrial environment.
  \item \textbf{System integration:} Contributed to the integration and maintenance of software services interacting with field equipment and operational charging infrastructure.
  \item \textbf{Needs clarification:} Translated operational issues, data-flow constraints and user feedback into technical tasks and corrective actions.
  \item \textbf{Production follow-up:} Analysed service behaviour, data exchanges, communication issues and anomalies affecting operational continuity.
  \item \textbf{Recette / validation:} Checked data consistency, interface continuity and service behaviour in production-like or operational conditions.
  \item \textbf{Issue management:} Investigated incidents involving software behaviour, operational data exchanges and hardware-connected systems.
  \item \textbf{Coordination:} Worked with technical stakeholders to follow priorities, clarify issues and support iterative delivery.
  \item \textbf{Data exchange:} Implemented and monitored FTP-based operational data exchange services.
  \item \textbf{Agile tooling:} Used Zenhub to support task tracking, backlog follow-up, iterative development and technical coordination.
\end{itemize}

\divider

\cvevent{\textbf{Senior R\&D Engineer - SI-like Workflows, Validation, Project Coordination \& Reporting}}{CNES, Laboratoire Lagrange, Observatoire de la Côte d'Azur}{2023 -- 2025}{Nice, France}

\textbf{Context:} Engineering activities for the LISA ESA/NASA space mission, involving complex software chains, model integration, simulation workflows, validation campaigns and international coordination.

\textbf{Objective:} Structure, coordinate and validate complex software/data workflows under strong consistency, traceability, performance and verification constraints.

\begin{itemize}
  \item \textbf{Project coordination:} Coordinated technical activities involving contributors, research engineers, technicians and Master's thesis interns.
  \item \textbf{Cadrage:} Structured development priorities, validation logic, expected outputs and technical acceptance criteria for complex analysis workflows.
  \item \textbf{Solution implementation:} Developed CATpyLISA, a Python platform for large-scale model integration, comparison, validation and technical review.
  \textcolor{DarkPastelRed}{\href{https://gitlab.oca.eu/gboileau/lisa_l3_tools}{\bf GitLab: CATpyLISA}}
  \item \textbf{V-cycle logic:} Supported activities aligned with specification, integration, verification, validation and technical evidence consolidation.
  \item \textbf{Recette / validation:} Defined comparison campaigns, consistency checks, acceptance logic and validation workflows for technical outputs.
  \item \textbf{Risk follow-up:} Identified deviations, inconsistencies, limiting factors and technical risks affecting result reliability or delivery quality.
  \item \textbf{Reporting:} Produced technical notes, validation summaries, analysis reports and review material for contributors and project stakeholders.
  \item \textbf{Documentation:} Used Confluence to support technical documentation, coordination notes, validation follow-up and knowledge sharing.
  \item \textbf{Stakeholder alignment:} Worked at the interface between scientific requirements, software development, data processing, modelling assumptions and mission-level objectives.
  \item \textbf{International environment:} Operated in a collaborative framework requiring technical English, structured communication and versioned workflows.
\end{itemize}

\divider

\cvevent{\textbf{Data Scientist - Technical Studies, Performance Analysis \& Project Reporting}}{Universiteit Antwerpen, Belgium}{2021 -- 2023}{Antwerp, Belgium}

\textbf{Context:} Data analysis and performance studies for precision instruments in a high-requirement international environment linked to gravitational-wave detector performance.

\textbf{Objective:} Identify, characterize and mitigate factors affecting system sensitivity using statistical, computational and validation-oriented methods.

\begin{itemize}
  \item \textbf{Technical studies:} Analysed instrumental and environmental effects impacting system sensitivity, stability and measurement quality.
  \item \textbf{Analysis workflows:} Developed Python-based pipelines for noise source identification, uncertainty estimation and statistical modelling.
  \item \textbf{Performance analysis:} Identified contributors to performance degradation and supported prioritisation of mitigation directions.
  \item \textbf{Data-driven decisions:} Analysed measurement and simulation outputs to identify abnormal behaviours and dominant contributors.
  \item \textbf{Documentation:} Produced reproducible and traceable technical analyses for international engineering and research stakeholders.
  \item \textbf{Coordination:} Communicated complex technical results through structured reporting, presentations and collaborative review.
\end{itemize}

\divider

\cvevent{\textbf{Junior R\&D Engineer - Simulation, Software Workflows \& Validation}}{ARTEMIS, Observatoire de la Côte d'Azur}{2018 -- 2021}{Nice, France}

\textbf{Context:} Engineering and analysis activities for gravitational-wave mission preparation, involving signal analysis, simulation, software workflows and noise characterization.

\textbf{Objective:} Develop robust analysis methods and validation workflows for complex signals and demanding system environments.

\begin{itemize}
  \item \textbf{Simulation workflows:} Built simulation approaches for complex signal environments, uncertainty studies and large-scale datasets.
  \item \textbf{Software validation:} Compared simulated outputs, model behaviours and analysis results to assess consistency and robustness.
  \item \textbf{Technical investigations:} Investigated noise sources through cross-sensor correlation analyses in diagnostic contexts.
  \item \textbf{Performance analysis:} Supported the identification of limiting factors affecting mission-level performance and system sensitivity.
  \item \textbf{Scientific computing:} Developed strong expertise in Python-based analysis, modelling, documentation and reproducible workflows.
\end{itemize}



\switchcolumn

\cvsection{Profile for Davidson}

\textbf{Technical project manager / systems engineer with strong experience in complex software workflows, technical coordination, validation, documentation, reporting and delivery support}. Background developed through international space and scientific collaborations, complemented by recent industrial software experience involving operational services, production follow-up and hardware/software interfaces.

For a \textbf{Chef de Projet SI} role, I bring a practical technical-management profile: analysing needs, coordinating technical and business-facing contributors, structuring deliverables, following risks and anomalies, supporting recette / validation, documenting decisions and helping users or stakeholders understand technical changes.

\begin{itemize}
  \item Experience coordinating technical activities across software, data, modelling, validation and operational constraints.
  \item Ability to translate needs, constraints and issues into structured actions, deliverables and follow-up.
  \item Experience with cadrage, specification logic, validation workflows, recette, technical reporting and documentation.
  \item Strong ability to analyse complex information, identify risks, propose pragmatic solutions and communicate clearly.
  \item Experience with Agile delivery contexts, iterative development, backlog follow-up and technical stakeholder coordination.
  \item Familiar with digital systems, applications, APIs, data flows, containers, Linux environments and software integration.
  \item Comfortable working in English and in multidisciplinary environments involving technical and non-technical stakeholders.
\end{itemize}

\cvsection{Fit for SI Project Management}

\begin{itemize}
  \item Relevant background for project phases: cadrage, design support, implementation follow-up, recette, validation and delivery support.
  \item Strong experience producing technical documentation, validation reports, progress summaries and decision-support material.
  \item Experience coordinating contributors, clarifying priorities, identifying risks and supporting corrective actions.
  \item Practical understanding of software architecture components: APIs, backend services, data flows, containers, Linux and operational monitoring.
  \item Comfortable with collaborative tools including Git/GitLab, GitHub, Confluence, Zenhub and JIRA-style issue tracking.
  \item Ability to support change through documentation, knowledge sharing, user guidance and structured communication.
\end{itemize}

\cvsection{Languages}

\cvskill{English}{4}
\divider
\cvskill{Spanish}{2}
\divider

\medskip

\cvsection{Education}

\cvevent{PhD in Physics}{Université Côte d'Azur}{2018 -- 2021}{Nice, France}

\divider

\cvevent{Selective Graduate Program in Fundamental Physics (Magistère)}{Université Paris-Saclay}{2016 -- 2018}{Orsay, France}

\divider

\cvevent{MSc in Astronomy and Astrophysics}{Observatoire de Paris / Université Paris-Saclay}{2017 -- 2018}{Paris, France}

\divider

\cvevent{BSc in Fundamental Physics}{Université Paris-Saclay}{2015 -- 2016}{Orsay, France}

\end{paracol}


\cvsection{Professional Training}

\cvevent{Supervised Machine Learning (Regression \& Classification)}{DeepLearning.AI - Andrew Ng}{2020}{Coursera}

\divider

\cvevent{Recherche reproductible : principes méthodologiques pour une science transparente}{FUN MOOC -- Inria}{2025}{France Université Numérique}

\divider

\cvevent{Continuous Professional Training -- Software, Data, AI and Systems}{OpenClassrooms}{2024 -- 2026}{}

\small{
Python, SQL, Machine Learning, AI, Linux, JavaScript, TypeScript, Angular, Node.js, Django, REST APIs, C/C++, web development, algorithms and software engineering fundamentals.
}


\cvsection{Projects}

\cvevent{CNES / LISA -- Technical Project Coordination, Software Workflows \& Validation}{ESA/NASA Space Mission - Large-scale International Collaboration}{2023 -- 2025}{Nice, France}

Coordination and structuring of software, data-processing and validation workflows for the LISA space mission within a large international engineering and scientific collaboration.

\textbf{Project management relevance:} Cadrage of technical activities, contributor coordination, validation planning, delivery follow-up, risk identification, reporting, documentation and stakeholder alignment.

\textbf{SI relevance:} Work involving complex software workflows, data-processing chains, GitLab-based development, reproducible environments, technical documentation, validation evidence and versioned deliverables.

\textbf{Scale:} Major international space mission with an estimated budget of approximately 2 billion EUR over 10 years.

\divider

\cvevent{Einstein Telescope / LIGO--Virgo--KAGRA -- Technical Studies \& Collaborative Project Workflows}{International Collaborations}{2021 -- 2023}{Antwerp, Belgium}

Contribution to collaborative developments in data analysis, software methods and technical investigations within large-scale international scientific and engineering collaborations.

\textbf{Project management relevance:} Coordination of technical studies, reporting, documentation, contribution to collaborative review, prioritisation of analysis tasks and communication with international stakeholders.

\textbf{System impact:} Studies on environmental and instrumental noise sources supporting the identification of performance limitations and design constraints for next-generation gravitational-wave observatories.

\cvsection{Strengths}

\vspace{-.3cm}

\cvsubsection{Project management \& delivery}

{\small
\cvtag{Project Management}
\cvtag{Project Coordination}
\cvtag{Cadrage}
\cvtag{Needs Analysis}
\cvtag{Requirements Formalisation}
\cvtag{Functional Analysis}
\cvtag{Technical Analysis}
\cvtag{Specifications}
\cvtag{Cahier des Charges}
\cvtag{Planning Follow-Up}
\cvtag{Risk Identification}
\cvtag{Action Plans}
\cvtag{Delivery Support}
\cvtag{Quality Follow-Up}
\cvtag{Budget Awareness}
\cvtag{Decision Support}
}

\divider\smallskip
\vspace{-.3cm}
\newpage

\cvsubsection{SI, recette \& deployment}

{\small
\cvtag{Information Systems}
\cvtag{SI Project Management}
\cvtag{Applications}
\cvtag{Digital Solutions}
\cvtag{Software Workflows}
\cvtag{APIs}
\cvtag{REST API}
\cvtag{Data Flows}
\cvtag{Backend Services}
\cvtag{Recette}
\cvtag{Validation}
\cvtag{Test Strategy}
\cvtag{Test Reports}
\cvtag{Deployment Support}
\cvtag{Production Follow-Up}
\cvtag{Mise en Production}
}

\divider\smallskip
\vspace{-.3cm}

\cvsubsection{Architecture, cloud \& technical environment}

{\small
\cvtag{Architecture SI Awareness}
\cvtag{Cloud Awareness}
\cvtag{Infrastructure Awareness}
\cvtag{Cybersecurity Awareness}
\cvtag{Linux}
\cvtag{Docker}
\cvtag{CI/CD}
\cvtag{Git / GitLab}
\cvtag{GitHub}
\cvtag{SQL}
\cvtag{Python}
\cvtag{TypeScript}
\cvtag{Angular}
\cvtag{Node.js}
\cvtag{Operational Monitoring}
\cvtag{OpenSearch}
}

\divider\smallskip
\vspace{-.3cm}

\cvsubsection{Tools, methods \& change management}

{\small
\cvtag{JIRA}
\cvtag{Confluence}
\cvtag{Zenhub}
\cvtag{Trello Awareness}
\cvtag{Agile}
\cvtag{Scrum Awareness}
\cvtag{ITIL Awareness}
\cvtag{Change Management}
\cvtag{User Support}
\cvtag{User Guidance}
\cvtag{Workshop Facilitation}
\cvtag{Stakeholder Alignment}
\cvtag{Cross-Functional Coordination}
\cvtag{Technical Documentation}
\cvtag{Progress Reporting}
\cvtag{Technical English}
}

\divider\smallskip
\vspace{-.3cm}

\cvsubsection{Office \& communication}

{\small
\cvtag{PowerPoint}
\cvtag{Excel}
\cvtag{Word}
\cvtag{MS Office}
\cvtag{Executive Presentations}
\cvtag{Reporting Dashboards}
\cvtag{Indicator Follow-Up}
\cvtag{Synthesis}
\cvtag{Analytical Thinking}
\cvtag{Problem Solving}
\cvtag{Pedagogy}
\cvtag{Autonomy}
\cvtag{Adaptability}
\cvtag{Rigor}
}
\end{document}